% Source: https://artofproblemsolving.com/community/c3h2795209_2021_tx_mmc__texas_mathworks_math_contest
% Stable identity: aops-topic-2795209-q06
The expression equals $2xy$, where $x$ is the sum of reciprocals of the first pair and $y$ that of the second pair. The sum $x+y=\dfrac17+\dfrac19+\dfrac1{11}+\dfrac1{13}$ is fixed, so $xy$ is largest when $x$ and $y$ are as close as possible.

The three pairings give
\[\{7,9\}\,\{11,13\}:\ \dfrac{16}{63},\dfrac{24}{143};\quad \{7,11\}\,\{9,13\}:\ \dfrac{18}{77},\dfrac{22}{117};\quad \{7,13\}\,\{9,11\}:\ \dfrac{20}{91},\dfrac{20}{99}.\]
The last pair is the closest (their difference is the smallest), so the maximum is
\[2\cdot\dfrac{20}{91}\cdot\dfrac{20}{99}=\dfrac{800}{9009}.\]
Since $9009=3^2\cdot7\cdot11\cdot13$ shares no factor with $800=2^5\cdot5^2$, the numerator is $\boxed{800}$.
